\section{GaitNet}
\label{sec:method}

GaitNet has two parts (Fig.~\ref{fig:overview}). A planner chooses, at
every tick, which step to take, if any, and an execution layer turns
that step into a command that the controller in use lands where and
when the planner meant. Table~\ref{tab:settings} lists the main
settings.

\subsection{Choosing One Step per Tick}
\label{subsec:method-step}

GaitNet chooses steps one at a time, and often enough to change its
plan while a leg is still in the air. At every tick it may start at
most one step for one leg, while the other legs keep their current
commands. A swing lasts several ticks, so GaitNet may start a second
swing before the first one ends. Two feasibility rules decide which
legs may step: a leg that is already swinging cannot be chosen, and a
new step needs at least \numberstringnum{\corlMinContactLegs} legs in
scheduled stance, where a leg that is waiting to lift off counts as
lifted. At most \numberstringnum{\corlMaxConcurrentSwingLegs} legs
therefore swing at once. Beyond these rules, GaitNet never chooses a
gait phase, leg pair, or contact sequence. It is \emph{greedy},
committing only to the next step, and the contact sequences it
produces are \emph{acyclic}, following neither a gait cycle nor a gait
library; the literature also calls them \emph{non-gaited} or
\emph{free-gait}. At each tick, it simply asks whether any available
step is better than doing nothing. If one is and another leg is
already in the air, the swings overlap; if the no-step option wins,
the robot waits.

The candidates are safe cells near each hip. A grid of terrain rays around the hip marks a cell as safe
when the ground under it is within reach below the hip, and as unsafe
when it is a void or the ray misses. The unsafe cells are then grown
by \numberstringnum{\corlMarginCm} cells, so that no remaining safe
cell lies within \corlMarginCm\,cm of a void; no other reach, edge, or
kinematic rule is applied. For each leg that may step, GaitNet draws a
fixed number of safe cells at random, without repeats, and masks the
empty slots when a leg has fewer. Together with the no-step option,
these candidates form the \emph{choices} of a tick, the same set in
training and at deployment.

\subsection{GaitNet Architecture}
\label{subsec:method-network}

GaitNet turns the robot state, the candidate set, and the recent motion
into one score and one swing time per choice (Fig.~\ref{fig:overview}a).

GaitNet has three inputs. The first is the tracking state
\begin{equation}
  \mathbf{x_g} =
  \begin{bmatrix}
    \mathbf p_{f,xy}^{b} &
    p_{b,z} &
    \mathbf v_b &
    \boldsymbol\omega_b &
    \mathbf v^{usr} &
    \mathbf c &
    \mathbf g_b
  \end{bmatrix}^T ,
  \label{eq:gaitnet-x}
\end{equation}
which holds the $xy$ positions of the four feet relative to the trunk
($\mathbf p_{f,xy}^{b}$), the base height in the world frame
($p_{b,z}$), the base linear and angular velocity ($\mathbf v_b$,
$\boldsymbol\omega_b$), the velocity command ($\mathbf v^{usr}$), the
controller's \emph{scheduled} contact rather than a sensed one
($\mathbf c$), and projected gravity ($\mathbf g_b$); all vectors
except $p_{b,z}$ are in the trunk frame. The second input is the
set of $N$ candidates and the no-step option $\varnothing$. Candidate
$i$ is described by a vector $\mathbf q_i$: a one-hot code for its leg
and its offset from the hip; the no-step option has its own code. A last entry, a cost from a separate network in
\cite{Sullivan2025}, is kept at zero, because zeroing it improved
survival there. The third input is the
history $\mathbf h$ of the last second of motion. Each of its ticks
holds the base velocities, each foot's position relative to its hip
from the joint encoders, the scheduled contact, and each leg's last
\emph{touchdown error}, which is where the foot was at its scheduled
touchdown minus the command. All inputs are available on a real robot.

GaitNet first combines these inputs. A state encoder $f_s$ and a
history encoder $f_h$ build a context vector
\begin{equation}
  \mathbf z = f_s(\mathbf x_g) + W_h\, f_h(\mathbf h),
  \label{eq:context}
\end{equation}
where $W_h$ is a linear projection. A candidate encoder $f_c$, shared by
all candidates, embeds each candidate as $\mathbf e_i = f_c(\mathbf q_i)$,
and the no-step option has a learned embedding $\mathbf e_\varnothing$
instead. Each choice then becomes a token
$\mathbf t_i = \mathrm{ReLU}(W_t [\mathbf z; \mathbf e_i])$ that
combines the shared context with that choice.

So that a choice can be judged
against the alternatives, a \emph{candidate mixer} $M$ lets the tokens
attend to each other. It is a small pre-norm transformer encoder
\cite{vaswani_attention_2017}: each layer updates every token as
$\mathbf t \leftarrow \mathbf t + \mathrm{MHA}(\mathrm{LN}(\mathbf t))$
and then $\mathbf t \leftarrow \mathbf t + \mathrm{FF}(\mathrm{LN}(\mathbf
t))$, with multi-head attention MHA, layer normalization LN, and a
feed-forward block FF. Empty slots are masked as keys, so no token
attends to them, while the no-step option is visible to all.

After the mixer, a readout $g$ maps each mixed token
$\tilde{\mathbf t}_i$ to a feature $\mathbf y_i = g(\tilde{\mathbf t}_i)$,
from which two linear heads give a score and a swing time,
\begin{equation}
  \alpha_i = \mathbf w_\alpha^\top \mathbf y_i, \qquad
  T_i = T_{\min} + (T_{\max} - T_{\min})\,
    \sigma(\mathbf w_T^\top \mathbf y_i),
  \label{eq:heads}
\end{equation}
where $\sigma$ is the logistic function and $[T_{\min}, T_{\max}]$ is
the allowed swing-time range.

The scores of the valid choices are the
logits of a categorical distribution $\pi(i) \propto \exp(\alpha_i)$;
empty slots get zero probability, and the no-step option is always
allowed. In training, GaitNet samples a choice $i$ and then a swing
time $T \sim \mathcal N(T_i, \sigma_T^2)$ with a small fixed
$\sigma_T$, and the log probability of the action is
$\log \pi(i) + \log \mathcal N(T; T_i, \sigma_T^2)$. At deployment it
takes the top-scoring choice and its mean swing time $T_i$, and issues
no step if the no-step option wins.

GaitNet starts training as a simpler network. The projection $W_h$, the attention output projections, and the last layer
of each feed-forward block start at zero. At initialization, the mixer
is therefore the identity and the context ignores the history, so every
choice is scored on its own from $\mathbf x_g$ and $\mathbf q_i$.
GaitNet then computes exactly the \emph{per-candidate network} of
\cite{Sullivan2025}, whose trunk corresponds to our token layer and readout. During training, GaitNet
learns to use the history and to compare candidates; afterwards, a
candidate's score depends on the whole candidate set, so GaitNet runs
on the same sampled set at deployment as in training.

\subsection{Inverse Footstep Model}
\label{subsec:method-idm}

The execution layer is needed because a controller does not always
place the foot where it is told, and each controller misses in its own
way. Our training controller keeps re-aiming each swing at the moving
hip, so the foot lands past the commanded point, and the faster the
robot walks, the farther past it lands. Without correction, we measure
a shift of +\corlOvershootSlowMin{} to +\corlOvershootSlowMax\,cm along
the walking direction at \corlSimSpeedSlow\,m/s and
+\corlOvershootFastMin{} to +\corlOvershootFastMax\,cm at
\corlSimSpeedFast\,m/s, against the \corlMarginCm\,cm margin that
GaitNet keeps around voids. OCS2, in contrast, fixes the target in the
world when the step is commanded.

We therefore fit one inverse footstep model (IDM) per controller. For
controller $c$, the IDM $I^c$ approximates the inverse of $F_c$ in
Eq.~\eqref{eq:forward}, $\mathbf u \approx I^c(\mathbf s, \ell, T,
\mathbf o)$. GaitNet's chosen candidate $\mathbf u^\star$ is a cell it
judged safe to land on, so we treat it as the desired landing and send
\begin{equation}
  \mathbf u = I^c(\mathbf s, \ell, T, \mathbf u^\star),
  \qquad\text{so that}\qquad
  F_c(\mathbf s, \ell, T, \mathbf u) \approx \mathbf u^\star .
  \label{eq:idm}
\end{equation}
GaitNet is the same for every controller; only $I^c$ changes. This
split into a planner and an inverse model follows Xie et al.\
\cite{xie_fada_2026}.

The IDM itself is a small network that predicts a correction to the desired
landing,
\begin{equation}
  I^c_\theta(\mathbf s, \ell, T, \mathbf o) =
  \mathbf o + \Delta^c_\theta(\mathbf h_{\text{last}}, \ell, T, \mathbf o),
  \label{eq:idm-model}
\end{equation}
where $\mathbf h_{\text{last}}$, the newest tick of GaitNet's history,
stands in for the robot state, and $\Delta^c_\theta$ is a small MLP.
The last layer of $\Delta^c_\theta$ starts at zero, so an untrained IDM
is the identity and sends GaitNet's choice unchanged.

We fit each IDM by hindsight relabelling. We log steps that its
controller actually executes, each as a tuple
$(\mathbf s, \ell, T, \mathbf u, \mathbf o)$. Whatever the command
$\mathbf u$ was aimed at, it is by definition a command that lands the
foot at $\mathbf o$, so the pair from $\mathbf o$ to $\mathbf u$ is a
correct example of the inverse map. We therefore fit
\begin{equation}
  \min_\theta \; \mathbb E\,
  \big\| \mathbf u - I^c_\theta(\mathbf s, \ell, T, \mathbf o) \big\|^2
  \label{eq:idm-loss}
\end{equation}
by supervised regression. For the training controller, the steps come
from two earlier policies without randomization, at up to 30\% voids
and across our speed range; a policy's own commands cover only the
footholds it prefers, so we add Gaussian noise to every command to
widen the data around them. For OCS2, we log the steps of the frozen
GaitNet, run on OCS2 with the training controller's IDM, on world seeds
61--110, fewer than 20{,}000 steps in all.

\subsection{Training GaitNet with the IDM}
\label{subsec:method-training}

We train GaitNet, as the actor, with Proximal Policy Optimization (PPO)
\cite{Schulman_proximal_2017} in Isaac Lab \cite{mittal_orbit_2023},
with many robots in parallel and the training controller's IDM, fixed,
between GaitNet and the controller. The critic is
a separate, smaller network that reads the state and the candidates,
but not the history, and shares no weights with the actor; it scores
each choice, and a small MLP maps these scores to one state value.

\begin{table}[t]
  \centering
  \caption{Main settings}
  \label{tab:settings}
  \footnotesize
  \setlength{\tabcolsep}{3pt}
  % source: gaitnet-docs/methodology.md 3-6, training-setup.md 5-10,
  % experiment-setup.md 2. Full settings: gaitnet-docs (supplementary).
  \begin{tabular*}{\columnwidth}{@{\extracolsep{\fill}}ll@{}}
    \toprule
    \thd{Item} & \thd{Value} \\
    \midrule
    \grouprow{2}{GaitNet}
    Tick; swing time & \corlControlTickMs\,ms;
      \corlSwingDurationMin--\corlSwingDurationMax\,s \\
    Candidates per leg; choices & \corlTrainCandidatesPerLeg;
      \corlTotalCandidates{} with no-step \\
    Candidate grid; void margin & $\pm\corlGridReachCm$\,cm at
      \corlScanResCm\,cm; \corlMarginCm\,cm \\
    History & \corlHistSeconds\,s (\corlHistTicks{} ticks) \\
    Candidate mixer & \corlMixerLayers{} layers, width \corlMixerDim,
      \corlMixerHeads{} heads \\
    Parameters & \corlParamsGaitNetK{}k (per-candidate:
      \corlParamsPerCandidateK{}k) \\
    \grouprow{2}{Inverse footstep model}
    Network; data & MLP, 2 $\times$ \corlIdmHidden;
      \corlIdmFootsteps{} steps \\
    Held-out error & \corlIdmRmsCm\,cm RMS (\corlIdmIdentityRmsCm\,cm
      uncorrected) \\
    \grouprow{2}{Training (Eq.~\eqref{eq:reward})}
    Robots; rollout; learning rate & \corlParallelEnvs;
      \corlRolloutSteps{} steps; $\corlLearningRate$ \\
    $w_a$; $w_s$; $w_f$ & \corlAliveWeight; \corlActionPenalty;
      \corlTerminationPenalty \\
    $\sigma_v$; $\sigma_\omega$; $\sigma_T$ & \corlTrackXyStd\,m/s;
      \corlTrackYawStd\,rad/s; \corlDurationExploreStd\,s \\
    Curriculum & \corlNumLevels{} levels, 0--\corlMaxVoidPct\% voids \\
    \grouprow{2}{OCS2}
    Step lead & \corlGzLead\,s \\
    \bottomrule
  \end{tabular*}
\end{table}


The training task is to walk over ground with voids. The terrain is a grid of cells, a
fraction of which, the \emph{void fraction}, are deep voids. The robot
must follow a forward velocity command until a time-out; it fails
earlier if the trunk tilts or drops too far, or if a foot drops below
the ground, which means a foot in a void. Commands are forward-only
because the deployed robot maps terrain with a forward-facing camera.
The reward per tick is
\begin{multline}
  r = \Delta t\,\Big[\,w_a\,\mathbb{I}_{\text{alive}}
  + e^{-\|\mathbf v^{usr}_{xy} - \mathbf v_{b,xy}\|^2/\sigma_v^2}
  + e^{-(\omega^{usr}_z - \omega_{b,z})^2/\sigma_\omega^2} \\
  - w_s\,\mathbb{I}_{\text{step}}
  - w_f\,\mathbb{I}_{\text{fail}}\Big],
  \label{eq:reward}
\end{multline}
where $\Delta t$ is the tick, superscript $usr$ marks the commanded
planar and yaw velocities, and subscript $b$ the matching base
velocities. There are no gait, clearance, posture, or smoothness
terms. The tracking widths $\sigma_v$ and $\sigma_\omega$ are narrow
on purpose: with Isaac Lab's default width, a robot standing still
would still earn most of the tracking reward at our speeds.
Observations are not normalized and carry no added noise, and there
are no pushes.

To make this task gradually harder, difficulty levels
raise the void fraction step by step. Isaac Lab's default curriculum
kept every robot on the lowest level at our speeds, so a robot instead
\emph{progresses} if it walks a set share of the commanded distance.
It moves up with probability $p_{\text{up}}$ after a progressing
time-out and down with probability $p_{\text{down}}$ after a failure
or a stalled time-out, so the share of failed or stalled episodes
settles near $p_{\text{up}}/(p_{\text{up}} + p_{\text{down}}) = 1/4$.
Each robot's controller is also redrawn at every reset, so GaitNet sees
a spread of tracking quality, swing heights, landing errors, and
lift-off delays.

The IDM is fitted first, offline (Sec.~\ref{subsec:method-idm}), and
then held fixed, so GaitNet and the IDM are not trained concurrently.
During GaitNet's training, every chosen candidate passes through the
training controller's IDM before the controller adds its randomized errors, although
the IDM itself was fitted without them. GaitNet still sees the raw
candidates, and its touchdown errors are measured against the IDM's
output, so its history reports how far the controller missed the
corrected command.

\subsection{Deployment on OCS2}
\label{subsec:method-deployment}

On OCS2, GaitNet replaces the gait scheduler of the
\texttt{legged\_control} stack \cite{leggedcontrol,
grandia_perceptive_2022} in Gazebo, where OCS2's NMPC \cite{ocs2}
plans the motion and a weighted WBC tracks it. A gait node asks a
GaitNet service for a decision every tick, also while two legs swing.
The service builds the same kind of candidate set, with a slightly
larger margin, rebuilds the history with the training code, and
applies the IDM fitted on OCS2's own steps (Sec.~\ref{subsec:method-idm}).

The IDM fixes where the foot lands; when a step starts is set by how
steps are delivered to OCS2, which also starts each step on its own
schedule. Each step is therefore written into OCS2's contact schedule
at absolute times on the NMPC clock, starting a fixed lead after the
decision, and is merged with the steps already planned. The body reference also stays a short distance ahead of the
body, and OCS2 runs with its stock task settings.

On OCS2, GaitNet and the controller get terrain either from
the world's exact geometry or from a simulated RealSense D435 depth
camera. For the camera, we build a \emph{depth-vote terrain}: a fine
grid in which each depth frame votes hit or miss for the cells it
sees. A cell is walkable while its fading votes are positive and it is
far enough from unwalkable cells, a distance that grows with how far
the robot has walked since it last saw the cell. The simulated depth images carry no noise. Odometry comes from OCS2's leg-based
state estimator with a \emph{slip gate}, which ignores feet whose
velocity disagrees with the estimate.
